\documentclass[11pt]{article}

\usepackage[utf8]{inputenc}
\usepackage[T1]{fontenc}
\usepackage{lmodern}
\usepackage{geometry}
\usepackage{amsmath,amssymb,amsfonts,amsthm,mathtools}
\usepackage{bm}
\usepackage{enumitem}
\usepackage{microtype}
\usepackage{hyperref}
\usepackage{titlesec}
\usepackage{booktabs}
\usepackage{array}
\usepackage{setspace}
\usepackage{mathrsfs}
\usepackage{physics}

\geometry{margin=1in}
\hypersetup{colorlinks=true, linkcolor=black, urlcolor=blue, citecolor=black}
\setlist[enumerate]{topsep=0.4em,itemsep=0.35em}
\setlist[itemize]{topsep=0.3em,itemsep=0.25em}
\titleformat{\section}{\large\bfseries}{\thesection.}{0.5em}{}
\titleformat{\subsection}{\normalsize\bfseries}{\thesubsection.}{0.5em}{}
\setstretch{1.06}

\newcommand{\Aether}{\AE{}ther}
\newcommand{\Aflow}{\AE{}ther-flow}
\newcommand{\TheoryName}{\Aflow{} Interpretation of Relativity}
\newcommand{\TheoryProgram}{\Aether{} / \Aflow{} framework}

\newcommand{\CompletedMetric}{g^{\sharp}}

\newcommand{\Car}{\mathrm{car}}
\newcommand{\Cur}{\mathrm{cur}}
\newcommand{\Hom}{\mathrm{hom}}

\newcommand{\ForSpace}[1]{\mathcal I_{#1}^{\mathrm{for}}}
\newcommand{\PrimShell}{\mathcal S_{\mathrm{prim}}}

\newcommand{\Reduction}[1]{\mathcal R_{#1}}
\newcommand{\CurrentCarrier}[1]{\mathfrak C_{#1}^{\mathrm{cur}}}
\newcommand{\EqCarrier}[1]{\mathfrak C_{#1}^{\mathrm{eq}}}
\newcommand{\CurrentExtract}[1]{\mathscr E_{#1}^{\mathrm{cur}}}
\newcommand{\EqExtract}[1]{\mathscr E_{#1}^{\mathrm{eq}}}
\newcommand{\PrimCurrent}[1]{\mathcal J_{#1}^{\mathrm{prim}}}
\newcommand{\PrimTheta}[1]{\Theta_{#1}^{\mathrm{prim}}}

\newcommand{\MetricOut}{\mathscr G_{\mathrm{out}}}
\newcommand{\MatterOut}{\mathscr T_{\mathrm{out}}}

\newcommand{\VisibleAffine}[1]{\widehat X_{#1}}
\newcommand{\DataSpace}[1]{\mathcal Y_{#1}^{\mathrm{obs}}}
\newcommand{\AmbientTarget}[1]{E_{#1}^{\mathrm{amb}}}

\newcommand{\EqnSpace}[1]{\mathcal U_{#1}^{\mathrm{eqn}}}
\newcommand{\EqnBody}[1]{\mathcal B_{#1}^{\mathrm{eqn}}}
\newcommand{\EqnData}[1]{\Lambda_{#1}^{\mathrm{eqn,dat}}}
\newcommand{\EqnShell}[1]{\Lambda_{#1}^{\mathrm{eqn,sh}}}
\newcommand{\EqnTarget}[1]{Q_{#1}^{\mathrm{eqn}}}
\newcommand{\EqnPair}[1]{\mathfrak b_{#1}^{\mathrm{eqn}}}
\newcommand{\EqnHidden}[1]{H_{#1}^{\mathrm{eqn}}}
\newcommand{\EqnZero}[1]{V_{#1}^{\mathrm{eqn},0}}

\newcommand{\ResSpace}[1]{\mathcal U_{#1}^{\mathrm{sup,res}}}
\newcommand{\ResBody}[1]{\mathcal B_{#1}^{\mathrm{sup,res}}}
\newcommand{\ResTarget}[1]{S_{#1}^{\mathrm{sup,res}}}
\newcommand{\ResPair}[1]{\mathfrak m_{#1}^{\mathrm{sup,res}}}
\newcommand{\ResSection}[1]{\Theta_{#1}^{\mathrm{sup,res}}}
\newcommand{\ResData}[1]{\Lambda_{#1}^{\mathrm{sup,dat}}}
\newcommand{\ResShellMap}[1]{\Lambda_{#1}^{\mathrm{sup,sh}}}
\newcommand{\ResShellSet}[1]{\mathcal Z_{#1}^{\mathrm{sup,res}}}
\newcommand{\SeedHidden}[1]{\mathcal H_{#1}^{\mathrm{sup}}}
\newcommand{\SeedHiddenBody}[1]{D_{#1}^{\mathrm{sup}}}
\newcommand{\HiddenChart}[1]{\Psi_{#1}^{\mathrm{hid}}}
\newcommand{\ZeroBody}[1]{\mathcal B_{#1}^{0,\mathrm{sup}}}

\newcommand{\SubSpace}[1]{\mathcal M_{#1}^{\mathrm{sub}}}
\newcommand{\SubBody}[1]{\mathcal B_{#1}^{\mathrm{sub}}}
\newcommand{\SubTarget}[1]{E_{#1}^{\mathrm{sub}}}
\newcommand{\SubPair}[1]{\mathfrak s_{#1}^{\mathrm{sub}}}
\newcommand{\SubShellSet}[1]{\mathcal Z_{#1}^{\mathrm{sub}}}

\newcommand{\LiftSpace}[1]{\mathcal L_{#1}^{\mathrm{lift}}}
\newcommand{\LiftBody}[1]{\mathcal B_{#1}^{\mathrm{lift}}}
\newcommand{\LiftProj}[1]{\Pi_{#1}^{\mathrm{lift}}}
\newcommand{\LiftSelect}[1]{\mathfrak s_{#1}^{\mathrm{lift}}}
\newcommand{\LiftSection}[1]{\Gamma_{#1}^{\mathrm{lift}}}
\newcommand{\LiftFunctional}[1]{\mathscr V_{#1}^{\mathrm{lift}}}
\newcommand{\LiftData}[1]{\Lambda_{#1}^{\mathrm{lift,dat}}}
\newcommand{\LiftShellMap}[1]{\Lambda_{#1}^{\mathrm{lift,sh}}}
\newcommand{\LiftSym}[1]{\mathbb G_{#1}^{\mathrm{lift}}}
\newcommand{\LiftTime}[1]{\vartheta_{#1}^{\mathrm{lift}}}
\newcommand{\LiftShellSet}[1]{\mathcal Z_{#1}^{\mathrm{lift}}}
\newcommand{\LiftChart}[1]{\Xi_{#1}^{\mathrm{lift}}}
\newcommand{\LiftSupport}[1]{\Theta_{#1}^{\mathrm{lift,sup}}}
\newcommand{\LiftSupportShell}[1]{\mathcal Z_{#1}^{\mathrm{lift,sup}}}
\newcommand{\LiftZeroCore}[1]{\mathcal Z_{#1}^{0,\mathrm{lift}}}
\newcommand{\LiftEqChart}[1]{\Psi_{#1}^{\mathrm{lift,sup}}}
\newcommand{\LiftZeroChart}[1]{\Psi_{#1}^{0,\mathrm{lift}}}
\newcommand{\LiftSplit}[1]{\Phi_{#1}^{\mathrm{lift}}}
\newcommand{\LiftCharge}[1]{\mathcal Q_{#1}^{\mathrm{lift}}}
\newcommand{\LiftResidual}[1]{\mathcal R_{#1}^{\mathrm{lift}}}
\newcommand{\LiftEmbed}[1]{\zeta_{#1}^{\mathrm{lift}}}
\newcommand{\LiftKernel}[1]{K_{#1}^{0,\mathrm{lift}}}
\newcommand{\LiftKernelCh}[1]{\widetilde K_{#1}^{0,\mathrm{lift}}}
\newcommand{\LiftStateComp}[1]{P_{#1}^{\mathrm{lift,co}}}

\newcommand{\PrecSpace}[1]{\mathcal P_{#1}^{\mathrm{prec}}}
\newcommand{\PrecBody}[1]{\mathcal B_{#1}^{\mathrm{prec}}}
\newcommand{\PrecProj}[1]{\Pi_{#1}^{\mathrm{prec}}}
\newcommand{\PrecLiftMin}[1]{\mathfrak f_{#1}^{\mathrm{prec}}}
\newcommand{\PrecSubMin}[1]{\mathfrak m_{#1}^{\mathrm{prec}}}
\newcommand{\PrecSection}[1]{\Gamma_{#1}^{\mathrm{prec}}}
\newcommand{\PrecFunctional}[1]{\mathscr W_{#1}^{\mathrm{prec}}}
\newcommand{\PrecData}[1]{\Lambda_{#1}^{\mathrm{prec,dat}}}
\newcommand{\PrecShellMap}[1]{\Lambda_{#1}^{\mathrm{prec,sh}}}
\newcommand{\PrecSym}[1]{\mathbb G_{#1}^{\mathrm{prec}}}
\newcommand{\PrecTime}[1]{\vartheta_{#1}^{\mathrm{prec}}}
\newcommand{\PrecShellSet}[1]{\mathcal Z_{#1}^{\mathrm{prec}}}
\newcommand{\PrecChart}[1]{\Xi_{#1}^{\mathrm{prec}}}
\newcommand{\PrecSupport}[1]{\Theta_{#1}^{\mathrm{prec,sup}}}
\newcommand{\PrecSupportShell}[1]{\mathcal Z_{#1}^{\mathrm{prec,sup}}}
\newcommand{\PrecZeroCore}[1]{\mathcal Z_{#1}^{0,\mathrm{prec}}}
\newcommand{\PrecEqChart}[1]{\Psi_{#1}^{\mathrm{prec,sup}}}
\newcommand{\PrecZeroChart}[1]{\Psi_{#1}^{0,\mathrm{prec}}}
\newcommand{\PrecSplit}[1]{\Phi_{#1}^{\mathrm{prec}}}
\newcommand{\PrecCharge}[1]{\mathcal Q_{#1}^{\mathrm{prec}}}
\newcommand{\PrecResidual}[1]{\mathcal R_{#1}^{\mathrm{prec}}}
\newcommand{\PrecEmbed}[1]{\zeta_{#1}^{\mathrm{prec}}}
\newcommand{\PrecKernel}[1]{K_{#1}^{0,\mathrm{prec}}}
\newcommand{\PrecKernelCh}[1]{\widetilde K_{#1}^{0,\mathrm{prec}}}
\newcommand{\PrecStateComp}[1]{P_{#1}^{\mathrm{prec,co}}}

\newcommand{\OrigSpace}[1]{\mathcal O_{#1}^{\mathrm{orig}}}
\newcommand{\OrigBody}[1]{\mathcal B_{#1}^{\mathrm{orig}}}
\newcommand{\OrigProj}[1]{\Pi_{#1}^{\mathrm{orig}}}
\newcommand{\OrigPrecMin}[1]{\mathfrak f_{#1}^{\mathrm{orig,prec}}}
\newcommand{\OrigLiftMin}[1]{\mathfrak f_{#1}^{\mathrm{orig,lift}}}
\newcommand{\OrigSubMin}[1]{\mathfrak m_{#1}^{\mathrm{orig}}}
\newcommand{\OrigSection}[1]{\Gamma_{#1}^{\mathrm{orig}}}
\newcommand{\OrigFunctional}[1]{\mathscr W_{#1}^{\mathrm{orig}}}
\newcommand{\OrigData}[1]{\Lambda_{#1}^{\mathrm{orig,dat}}}
\newcommand{\OrigShellMap}[1]{\Lambda_{#1}^{\mathrm{orig,sh}}}
\newcommand{\OrigSym}[1]{\mathbb G_{#1}^{\mathrm{orig}}}
\newcommand{\OrigTime}[1]{\vartheta_{#1}^{\mathrm{orig}}}
\newcommand{\OrigShellSet}[1]{\mathcal Z_{#1}^{\mathrm{orig}}}
\newcommand{\OrigChart}[1]{\Xi_{#1}^{\mathrm{orig}}}
\newcommand{\OrigSupport}[1]{\Theta_{#1}^{\mathrm{orig,sup}}}
\newcommand{\OrigSupportShell}[1]{\mathcal Z_{#1}^{\mathrm{orig,sup}}}
\newcommand{\OrigZeroCore}[1]{\mathcal Z_{#1}^{0,\mathrm{orig}}}
\newcommand{\OrigEqChart}[1]{\Psi_{#1}^{\mathrm{orig,sup}}}
\newcommand{\OrigZeroChart}[1]{\Psi_{#1}^{0,\mathrm{orig}}}
\newcommand{\OrigSplit}[1]{\Phi_{#1}^{\mathrm{orig}}}
\newcommand{\OrigCharge}[1]{\mathcal Q_{#1}^{\mathrm{orig}}}
\newcommand{\OrigResidual}[1]{\mathcal R_{#1}^{\mathrm{orig}}}
\newcommand{\OrigEmbed}[1]{\zeta_{#1}^{\mathrm{orig}}}
\newcommand{\OrigKernel}[1]{K_{#1}^{0,\mathrm{orig}}}
\newcommand{\OrigKernelCh}[1]{\widetilde K_{#1}^{0,\mathrm{orig}}}
\newcommand{\OrigStateComp}[1]{P_{#1}^{\mathrm{orig,co}}}

\newtheorem{definition}{Definition}
\newtheorem{lemma}{Lemma}
\newtheorem{proposition}{Proposition}
\newtheorem{remark}{Remark}
\newtheorem{corollary}{Corollary}
\newtheorem{theorem}{Theorem}

\title{\textbf{The \TheoryName{}}\\[0.4em]
\large Primitive-Reservoir Observer-Reduction\\
\large Lift-Generating Substrate Precursor Dynamics\\
\large from Precursor-Generating Substrate Origin Dynamics}
\author{}
\date{}

\begin{document}

\maketitle

\begin{abstract}
The current primitive-reservoir observer-reduction frontier has already moved the honest burden below the shell-generating substrate lift layer. The remaining live question is now one layer lower again. The lift-generating substrate precursor package is no longer the right place to stop if the project goal is a disciplined deeper derivation of the same benchmark one-metric package. What therefore remains open is not another shell-, lift-, support-, or reservoir-level restatement, not a reopening of stationary reduction, stationary Hessians, spectral generators, or selector mechanics, and not any change to the one operative metric or the ordinary matter law. The remaining honest burden is to derive or justify the lift-generating substrate precursor dynamics themselves from still more primitive substrate structure, or else to prove a sharper inevitability theorem at that deeper level.

The present note takes the deeper origin-side route. For each sector it introduces precursor-generating substrate origin dynamics: a compact convex origin body over the precursor state space, a smooth origin restriction section with unique origin minimizers on fixed-precursor, fixed-lift, and fixed-substrate fibers, affine origin observer-data and shell-response readouts factoring through the precursor readouts, symmetry and time reversal, an origin exact shell projecting diffeomorphically to the exact precursor shell, an origin shell chart to the same reservoir body, an origin shell-internal support recorder, an origin support shell with zero-core / hidden-seed decomposition, an origin support-shell chart to the same equation variables, an origin hidden charge, an observer-compatible exact origin zero-response realization, fixed-data solvability, and coercive control on origin data-invisible directions. The lift-generating substrate precursor dynamics are then recovered canonically by projected exact-shell transport from that deeper origin package.

Because the induced precursor package is exactly the one already used by the immediately preceding lift-generation theorem, the downstream lift, shell-restriction, and reservoir-facing consequences remain available unchanged. The same benchmark one-metric package remains fixed throughout: the operative metric is still exactly \(\CompletedMetric\), the ordinary matter route is unchanged, there is no second operative metric, and the low-energy matter law is not enlarged. The scope remains disciplined. This note is an origin-side theorem inside the active derivational program of \TheoryName{}, not a license for stronger promotion language. The next honest burden is deeper again: derive or justify the precursor-generating substrate origin dynamics themselves from still more primitive substrate structure, or else prove a sharper inevitability theorem there.
\end{abstract}

\section{Introduction}

The active derivational program of \TheoryName{} now has a sharper next task than another same-output relay. The immediately preceding frontier theorem already showed that the shell-generating substrate lift package is no longer the disciplined place to stop on the active primitive-reservoir observer-reduction line. Once lift-generating substrate precursor dynamics are granted, the entire lift package is already the canonical projected exact-shell output of that deeper precursor class. That result matters precisely because it moves the honest burden below the lift layer. The next manuscript therefore cannot honestly be another shell-, lift-, support-, or reservoir-level reformulation, another reopening of stationary reduction, stationary Hessians, spectral generators, or selector layers, or any change to the one operative metric and the ordinary matter route. The next manuscript has to address the precursor dynamics themselves.

There are two disciplined ways to do that. One can descend by one more explicit origin layer and derive the precursor package as the exact projected consequence of still deeper substrate data. Or one can prove that, inside an even sharper deeper class, the precursor package is itself forced. The present note takes the first route. It treats the current precursor package as the recorded target, introduces a more primitive origin layer below it, and proves that the entire lift-generating substrate precursor package arises there as a projected exact-shell transport. In that sense the precursor layer is no longer left as a free insertion on the active line.

That is the correct scope for the next main-line step. The point is not to change the benchmark exact-relativistic theory statement, and it is not to reopen already-settled line-routing decisions. The point is to remove the precursor package itself as the next unexplained insertion on the active line. The positive gain is therefore narrower than a completed first-principles substrate derivation and stronger than another same-output restatement of already recorded lower layers.

\paragraph{Framework and claim boundary.}
\TheoryProgram{} denotes the broader \Aether{} / \Aflow{} framework built on the claims that the \Aether{} is the underlying four-dimensional substrate of reality, the \Aflow{} is its intrinsic ordered motion, observed three-dimensional space is the local experiential slice of that deeper substrate, S-time is the experienced order of change arising from matter, light, and the \Aflow{}, and observed expansion is the three-dimensional appearance of deeper four-dimensional ordered motion. Within that framework, \TheoryName{} denotes the exact-closure theory in which the effective gravitational dynamics are adopted to be exactly Einsteinian with universal matter coupling. In the active sequence, \emph{adoption} means use of that established relativistic dynamics without claiming substrate derivation, while \emph{derivation} is reserved for a first-principles recovery from explicit substrate variables.

The present note stays strictly inside that boundary. It does not alter the benchmark exact-closure package. It does not introduce a second operative metric or a new low-energy matter law. It only strengthens the deeper origin-side explanation of why the lift-generating substrate precursor dynamics used by the previous theorem should exist.

\section{Fixed Primitive Continuation Data}

\begin{definition}[Recorded primitive continuation data]
\label{def:fixed}
Fix the following continuation data from the active primitive-reservoir observer-reduction line.
\begin{enumerate}
    \item For each sector label \(\sigma\in\{\Car,\Cur,\Hom\}\), a primitive forcing shell
    \[
    \ForSpace{\sigma},
    \]
    a visible reduction map
    \[
    \Reduction{\sigma}:\ForSpace{\sigma}\to X_{\sigma},
    \]
    and shell-level current and equilibrium carriers
    \[
    \CurrentCarrier{\sigma}:\ForSpace{\sigma}\to Z_{\sigma}^{\mathrm{cur}},
    \qquad
    \EqCarrier{\sigma}:\ForSpace{\sigma}\to Z_{\sigma}^{\mathrm{eq}}.
    \]
    \item The product shell
    \[
    \PrimShell
    \equiv
    \ForSpace{\Car}\times \ForSpace{\Cur}\times \ForSpace{\Hom}.
    \]
    \item Fixed shell extractors
    \[
    \CurrentExtract{\sigma}:Z_{\sigma}^{\mathrm{cur}}\to Y_{\sigma}^{\mathrm{cur}},
    \qquad
    \EqExtract{\sigma}:Z_{\sigma}^{\mathrm{eq}}\to Y_{\sigma}^{\mathrm{eq}},
    \]
    together with the recorded shell composites
    \begin{equation}
    \PrimCurrent{\sigma}
    =
    \CurrentExtract{\sigma}\circ \CurrentCarrier{\sigma},
    \qquad
    \PrimTheta{\sigma}
    =
    \EqExtract{\sigma}\circ \EqCarrier{\sigma}.
    \label{eq:primcomposites}
    \end{equation}
    \item The recorded metric and ordinary-matter outputs
    \[
    \MetricOut:\PrimShell\to\{\text{observer metrics}\},
    \qquad
    \MatterOut:\PrimShell\times\Psi\to\{\text{observer stress tensors}\},
    \]
    satisfying
    \begin{equation}
    \MetricOut(\mathbf w)=\CompletedMetric,
    \qquad
    \MatterOut(\mathbf w,\psi)
    =
    T_{\mu\nu}^{\mathrm{matter}}[\CompletedMetric,\psi]
    \label{eq:fixedoutputs}
    \end{equation}
    for every \(\mathbf w\in\PrimShell\) and every matter configuration \(\psi\in\Psi\).
\end{enumerate}
\end{definition}

\begin{remark}
Definition \ref{def:fixed} is not reopened here. The primitive continuation data, the one operative metric, and the ordinary matter route remain fixed. The present note only addresses the deeper origin-side status of the lift-generating substrate precursor package that now sits at the frontier of the active line.
\end{remark}

\section{Recorded Lift-Generating Precursor Target}

\begin{definition}[Lift-generating substrate precursor dynamics]
\label{def:prec}
Fix \(\sigma\in\{\Car,\Cur,\Hom\}\). A \emph{lift-generating substrate precursor dynamics package} for sector \(\sigma\) consists of the following data.
\begin{enumerate}
    \item A finite-dimensional Euclidean precursor state space
    \[
    \PrecSpace{\sigma},
    \]
    a compact convex precursor body
    \[
    \PrecBody{\sigma}\subset\PrecSpace{\sigma}
    \]
    with nonempty interior, an affine surjection
    \[
    \PrecProj{\sigma}:\PrecSpace{\sigma}\to\LiftSpace{\sigma},
    \]
    and the projected lift body
    \[
    \LiftBody{\sigma}
    \equiv
    \PrecProj{\sigma}\bigl(\PrecBody{\sigma}\bigr)
    \subset\LiftSpace{\sigma},
    \]
    assumed compact convex with nonempty interior. There is an affine surjection
    \[
    \LiftProj{\sigma}:\LiftSpace{\sigma}\to\SubSpace{\sigma},
    \]
    and the projected substrate body
    \[
    \SubBody{\sigma}
    \equiv
    \LiftProj{\sigma}\bigl(\LiftBody{\sigma}\bigr)
    \subset\SubSpace{\sigma},
    \]
    also assumed compact convex with nonempty interior.

    There is a smooth precursor lift-fiber minimizer
    \[
    \PrecLiftMin{\sigma}:\LiftBody{\sigma}\to\PrecBody{\sigma}
    \]
    satisfying
    \[
    \PrecProj{\sigma}\circ\PrecLiftMin{\sigma}
    =
    \mathrm{id}_{\LiftBody{\sigma}},
    \]
    and a smooth precursor substrate-fiber minimizer
    \[
    \PrecSubMin{\sigma}:\SubBody{\sigma}\to\PrecBody{\sigma}
    \]
    satisfying
    \[
    \LiftProj{\sigma}\circ\PrecProj{\sigma}\circ\PrecSubMin{\sigma}
    =
    \mathrm{id}_{\SubBody{\sigma}}.
    \]
    Define
    \begin{equation}
    \LiftSelect{\sigma}
    \equiv
    \PrecProj{\sigma}\circ\PrecSubMin{\sigma}.
    \label{eq:inducedliftselect}
    \end{equation}
    Assume
    \begin{equation}
    \PrecLiftMin{\sigma}\bigl(\LiftSelect{\sigma}(m)\bigr)
    =
    \PrecSubMin{\sigma}(m)
    \qquad
    \text{for every }m\in\SubBody{\sigma}.
    \label{eq:compatmin}
    \end{equation}
    \item The same finite-dimensional Euclidean substrate restriction target
    \[
    \SubTarget{\sigma},
    \]
    the same positive-definite symmetric substrate pairing
    \[
    \SubPair{\sigma}:
    \SubTarget{\sigma}\times\SubTarget{\sigma}\to\mathbb R,
    \]
    and a smooth precursor restriction section
    \[
    \PrecSection{\sigma}:\PrecSpace{\sigma}\to\SubTarget{\sigma}.
    \]
    Define the precursor restriction functional
    \begin{equation}
    \PrecFunctional{\sigma}(p)
    \equiv
    \SubPair{\sigma}\bigl(\PrecSection{\sigma}(p),\PrecSection{\sigma}(p)\bigr),
    \label{eq:precfunctional}
    \end{equation}
    together with the exact precursor shell
    \begin{equation}
    \PrecShellSet{\sigma}
    \equiv
    \bigl(\PrecSection{\sigma}\bigr)^{-1}(0)\cap\PrecBody{\sigma}.
    \label{eq:precshell}
    \end{equation}
    Define the induced lifted restriction section
    \begin{equation}
    \LiftSection{\sigma}
    \equiv
    \PrecSection{\sigma}\circ\PrecLiftMin{\sigma},
    \label{eq:inducedliftsection}
    \end{equation}
    and the induced lifted restriction functional
    \begin{equation}
    \LiftFunctional{\sigma}(\ell)
    \equiv
    \SubPair{\sigma}\bigl(\LiftSection{\sigma}(\ell),\LiftSection{\sigma}(\ell)\bigr)
    =
    \PrecFunctional{\sigma}\bigl(\PrecLiftMin{\sigma}(\ell)\bigr).
    \label{eq:inducedliftfunctional}
    \end{equation}
    Assume that, for each \(\ell\in\LiftBody{\sigma}\), \(\PrecLiftMin{\sigma}(\ell)\) is the unique minimizer of \(\PrecFunctional{\sigma}\) on the fiber
    \[
    \bigl(\PrecProj{\sigma}\bigr)^{-1}(\ell)\cap\PrecBody{\sigma},
    \]
    and, for each \(m\in\SubBody{\sigma}\), \(\PrecSubMin{\sigma}(m)\) is the unique minimizer of \(\PrecFunctional{\sigma}\) on the composite fiber
    \[
    \bigl(\LiftProj{\sigma}\circ\PrecProj{\sigma}\bigr)^{-1}(m)\cap\PrecBody{\sigma}.
    \]
    Assume that \(\LiftSelect{\sigma}(m)\) is therefore the unique minimizer of \(\LiftFunctional{\sigma}\) on
    \[
    \bigl(\LiftProj{\sigma}\bigr)^{-1}(m)\cap\LiftBody{\sigma},
    \]
    and that there exists \(\mu_{\sigma}^{\mathrm{fib}}>0\) such that
    \begin{equation}
    \LiftFunctional{\sigma}(\ell)
    \ge
    \LiftFunctional{\sigma}\bigl(\LiftSelect{\sigma}(m)\bigr)
    +
    \mu_{\sigma}^{\mathrm{fib}}\,
    \operatorname{dist}\bigl(\ell,\LiftSelect{\sigma}(m)\bigr)^{2}
    \label{eq:precfibergap}
    \end{equation}
    for every \(m\in\SubBody{\sigma}\) and every \(\ell\in\bigl(\LiftProj{\sigma}\bigr)^{-1}(m)\cap\LiftBody{\sigma}\). Assume there exists \(\mu_{\sigma}^{\mathrm{sub}}>0\) such that
    \begin{equation}
    \LiftFunctional{\sigma}\bigl(\LiftSelect{\sigma}(m)\bigr)
    \ge
    \mu_{\sigma}^{\mathrm{sub}}\,
    \operatorname{dist}\bigl(m,\LiftProj{\sigma}(\PrecProj{\sigma}(\PrecShellSet{\sigma}))\bigr)^{2}
    \qquad
    \text{for all }m\in\SubBody{\sigma}.
    \label{eq:precprojectedgap}
    \end{equation}
    \item Affine lifted observer-data and shell-response readouts
    \[
    \LiftData{\sigma}:\LiftSpace{\sigma}\to\DataSpace{\sigma},
    \qquad
    \LiftShellMap{\sigma}:\LiftSpace{\sigma}\to\AmbientTarget{\sigma},
    \]
    and affine precursor observer-data and shell-response readouts
    \[
    \PrecData{\sigma}:\PrecSpace{\sigma}\to\DataSpace{\sigma},
    \qquad
    \PrecShellMap{\sigma}:\PrecSpace{\sigma}\to\AmbientTarget{\sigma},
    \]
    satisfying
    \begin{equation}
    \PrecData{\sigma}
    =
    \LiftData{\sigma}\circ\PrecProj{\sigma},
    \qquad
    \PrecShellMap{\sigma}
    =
    \LiftShellMap{\sigma}\circ\PrecProj{\sigma}.
    \label{eq:precfactor}
    \end{equation}
    There is also a finite-dimensional Euclidean support target
    \[
    \ResTarget{\sigma},
    \]
    a positive-definite symmetric support pairing
    \[
    \ResPair{\sigma}:
    \ResTarget{\sigma}\times\ResTarget{\sigma}\to\mathbb R,
    \]
    a finite-dimensional real hidden target
    \[
    \EqnTarget{\sigma},
    \]
    and a positive-definite symmetric hidden pairing
    \[
    \EqnPair{\sigma}:
    \EqnTarget{\sigma}\times\EqnTarget{\sigma}\to\mathbb R.
    \]
    There is a compact symmetry group
    \[
    \PrecSym{\sigma}
    \]
    and an involution
    \[
    \PrecTime{\sigma},
    \]
    acting orthogonally on \(\PrecSpace{\sigma}\), \(\LiftSpace{\sigma}\), \(\SubSpace{\sigma}\), \(\DataSpace{\sigma}\), \(\AmbientTarget{\sigma}\), and \(\ResTarget{\sigma}\), and by \(\EqnPair{\sigma}\)-isometries on \(\EqnTarget{\sigma}\), preserving \(\PrecBody{\sigma}\), \(\LiftBody{\sigma}\), \(\PrecProj{\sigma}\), \(\LiftProj{\sigma}\), \(\PrecLiftMin{\sigma}\), \(\PrecSubMin{\sigma}\), \(\PrecSection{\sigma}\), \(\LiftData{\sigma}\), \(\LiftShellMap{\sigma}\), \(\PrecData{\sigma}\), and \(\PrecShellMap{\sigma}\).
    \item The restriction of \(\PrecProj{\sigma}\) to \(\PrecShellSet{\sigma}\) is a diffeomorphism onto its image
    \[
    \LiftShellSet{\sigma}
    \equiv
    \PrecProj{\sigma}\bigl(\PrecShellSet{\sigma}\bigr)
    \subset\LiftBody{\sigma}.
    \]
    There is a Euclidean precursor shell chart
    \[
    \PrecChart{\sigma}:\PrecShellSet{\sigma}\to\ResSpace{\sigma}
    \]
    whose shell-body image
    \[
    \ResBody{\sigma}
    \equiv
    \PrecChart{\sigma}\bigl(\PrecShellSet{\sigma}\bigr)
    \]
    is compact convex with nonempty interior. There exist affine maps
    \[
    \ResData{\sigma}:\ResSpace{\sigma}\to\DataSpace{\sigma},
    \qquad
    \ResShellMap{\sigma}:\ResSpace{\sigma}\to\AmbientTarget{\sigma},
    \]
    and a smooth support section
    \[
    \ResSection{\sigma}:\ResSpace{\sigma}\to\ResTarget{\sigma}
    \]
    such that
    \begin{equation}
    \ResData{\sigma}\circ\PrecChart{\sigma}
    =
    \PrecData{\sigma}\big|_{\PrecShellSet{\sigma}},
    \qquad
    \ResShellMap{\sigma}\circ\PrecChart{\sigma}
    =
    \PrecShellMap{\sigma}\big|_{\PrecShellSet{\sigma}}.
    \label{eq:precchartreadouts}
    \end{equation}
    There is a smooth precursor shell-internal support recorder
    \[
    \PrecSupport{\sigma}:\PrecShellSet{\sigma}\to\ResTarget{\sigma}
    \]
    satisfying
    \begin{equation}
    \ResSection{\sigma}\circ\PrecChart{\sigma}
    =
    \PrecSupport{\sigma}.
    \label{eq:precsupporttransport}
    \end{equation}
    Define the exact precursor support shell
    \begin{equation}
    \PrecSupportShell{\sigma}
    \equiv
    \bigl(\PrecSupport{\sigma}\bigr)^{-1}(0)\cap\PrecShellSet{\sigma},
    \label{eq:precsupportshell}
    \end{equation}
    and assume
    \begin{equation}
    \ResShellSet{\sigma}
    \equiv
    \bigl(\ResSection{\sigma}\bigr)^{-1}(0)\cap\ResBody{\sigma}
    =
    \PrecChart{\sigma}\bigl(\PrecSupportShell{\sigma}\bigr).
    \label{eq:prectransportedsupportshell}
    \end{equation}
    Assume there exists \(\nu_{\sigma}^{\mathrm{sup}}>0\) such that
    \begin{equation}
    \ResPair{\sigma}\bigl(\ResSection{\sigma}(u),\ResSection{\sigma}(u)\bigr)
    \ge
    \nu_{\sigma}^{\mathrm{sup}}\,
    \operatorname{dist}\bigl(u,\ResShellSet{\sigma}\bigr)^{2}
    \qquad
    \text{for all }u\in\ResBody{\sigma}.
    \label{eq:precnormal}
    \end{equation}
    Assume the restriction of \(\PrecProj{\sigma}\) to \(\PrecSupportShell{\sigma}\) is a diffeomorphism onto its image
    \[
    \LiftSupportShell{\sigma}
    \equiv
    \PrecProj{\sigma}\bigl(\PrecSupportShell{\sigma}\bigr).
    \]
    The transported symmetry and time-reversal actions act linearly on \(\ResSpace{\sigma}\), preserving \(\ResBody{\sigma}\), \(\ResShellSet{\sigma}\), \(\ResSection{\sigma}\), \(\ResData{\sigma}\), and \(\ResShellMap{\sigma}\).
    \item A closed embedded precursor support zero core
    \[
    \PrecZeroCore{\sigma}\subset\PrecSupportShell{\sigma},
    \]
    a finite-dimensional Euclidean hidden-seed space
    \[
    \SeedHidden{\sigma},
    \]
    and a compact convex hidden-seed body
    \[
    \SeedHiddenBody{\sigma}\subset\SeedHidden{\sigma}
    \]
    with
    \[
    0\in\operatorname{int}\SeedHiddenBody{\sigma}.
    \]
    There is a Euclidean precursor support-shell chart
    \[
    \PrecEqChart{\sigma}:\PrecSupportShell{\sigma}\to\EqnSpace{\sigma},
    \]
    an affine chart on the precursor support zero core
    \[
    \PrecZeroChart{\sigma}:\PrecZeroCore{\sigma}\to\EqnZero{\sigma},
    \]
    and a linear isometric embedding
    \[
    \HiddenChart{\sigma}:\SeedHidden{\sigma}\to\EqnHidden{\sigma}
    \]
    such that
    \[
    \EqnSpace{\sigma}
    =
    \EqnZero{\sigma}\oplus\EqnHidden{\sigma}
    \]
    is an orthogonal direct sum, the set
    \[
    \ZeroBody{\sigma}
    \equiv
    \PrecZeroChart{\sigma}\bigl(\PrecZeroCore{\sigma}\bigr)
    \]
    is compact convex with nonempty interior in \(\EqnZero{\sigma}\), the shell image
    \[
    \EqnBody{\sigma}
    \equiv
    \PrecEqChart{\sigma}\bigl(\PrecSupportShell{\sigma}\bigr)
    \]
    has nonempty interior in \(\EqnSpace{\sigma}\), and there is a diffeomorphism
    \begin{equation}
    \PrecSplit{\sigma}:
    \PrecZeroCore{\sigma}\times\SeedHiddenBody{\sigma}
    \xrightarrow{\cong}
    \PrecSupportShell{\sigma}
    \label{eq:precsplit}
    \end{equation}
    satisfying
    \begin{equation}
    \PrecEqChart{\sigma}\bigl(\PrecSplit{\sigma}(z,\eta)\bigr)
    =
    \PrecZeroChart{\sigma}(z)+\HiddenChart{\sigma}(\eta).
    \label{eq:precchartsplit}
    \end{equation}
    The precursor shell readouts become affine in the precursor support-shell chart: there exist affine maps
    \[
    \EqnData{\sigma}:\EqnSpace{\sigma}\to\DataSpace{\sigma},
    \qquad
    \EqnShell{\sigma}:\EqnSpace{\sigma}\to\AmbientTarget{\sigma},
    \]
    such that
    \begin{equation}
    \EqnData{\sigma}\circ\PrecEqChart{\sigma}
    =
    \PrecData{\sigma}\big|_{\PrecSupportShell{\sigma}},
    \qquad
    \EqnShell{\sigma}\circ\PrecEqChart{\sigma}
    =
    \PrecShellMap{\sigma}\big|_{\PrecSupportShell{\sigma}}.
    \label{eq:preceqnchartreadouts}
    \end{equation}
    Assume the restriction of \(\PrecProj{\sigma}\) to \(\PrecZeroCore{\sigma}\) is a diffeomorphism onto its image
    \[
    \LiftZeroCore{\sigma}
    \equiv
    \PrecProj{\sigma}\bigl(\PrecZeroCore{\sigma}\bigr).
    \]
    The transported symmetry and time-reversal actions act linearly on \(\EqnSpace{\sigma}\), preserving \(\EqnZero{\sigma}\), \(\EqnHidden{\sigma}\), and \(\EqnBody{\sigma}\).
    \item A linear hidden charge map
    \[
    \PrecCharge{\sigma}:\SeedHidden{\sigma}\to\EqnTarget{\sigma}
    \]
    such that there exists \(\lambda_{\sigma}^{\mathrm{eqn}}>0\) with
    \begin{equation}
    \EqnPair{\sigma}\bigl(\PrecCharge{\sigma}\eta,\PrecCharge{\sigma}\eta\bigr)
    \ge
    \lambda_{\sigma}^{\mathrm{eqn}}\,
    \|\eta\|^{2}
    \qquad
    \text{for all }\eta\in\SeedHidden{\sigma}.
    \label{eq:precchargegap}
    \end{equation}
    Define the precursor support-shell hidden recorder by
    \begin{equation}
    \PrecResidual{\sigma}\bigl(\PrecSplit{\sigma}(z,\eta)\bigr)
    \equiv
    \PrecCharge{\sigma}(\eta).
    \label{eq:precresidual}
    \end{equation}
    \item A smooth embedding
    \[
    \jmath_{\sigma}:X_{\sigma}\hookrightarrow\VisibleAffine{\sigma}
    \]
    and a smooth observer-compatible exact precursor zero-response realization
    \[
    \PrecEmbed{\sigma}:\ForSpace{\sigma}\hookrightarrow\PrecZeroCore{\sigma}
    \]
    whose image is exactly
    \[
    \PrecZeroCore{\sigma}
    \cap
    \bigl(\PrecShellMap{\sigma}\bigr)^{-1}(0),
    \]
    and such that
    \begin{equation}
    \PrecData{\sigma}\bigl(\PrecEmbed{\sigma}(w)\bigr)
    =
    \bigl(
    \jmath_{\sigma}(\Reduction{\sigma}(w)),
    \CurrentCarrier{\sigma}(w),
    \EqCarrier{\sigma}(w)
    \bigr)
    \label{eq:preccompat}
    \end{equation}
    for every \(w\in\ForSpace{\sigma}\).
    \item Fixed-data solvability on the precursor support zero core:
    \begin{equation}
    \PrecData{\sigma}\bigl(\PrecZeroCore{\sigma}\bigr)
    =
    \PrecData{\sigma}\Bigl(
    \PrecZeroCore{\sigma}\cap
    \bigl(\PrecShellMap{\sigma}\bigr)^{-1}(0)
    \Bigr).
    \label{eq:precsolvability}
    \end{equation}
    \item The inert precursor support-zero-core kernel
    \[
    \PrecKernel{\sigma}
    \equiv
    \ker\bigl(D\PrecData{\sigma}\big|_{\PrecZeroCore{\sigma}}\bigr)
    \cap
    \ker\bigl(D\PrecShellMap{\sigma}\big|_{\PrecZeroCore{\sigma}}\bigr)
    \]
    and the orthogonal projector
    \[
    \PrecStateComp{\sigma}:
    \EqnZero{\sigma}\to
    \bigl(\PrecKernelCh{\sigma}\bigr)^{\perp}\cap\EqnZero{\sigma},
    \qquad
    \PrecKernelCh{\sigma}
    \equiv
    D\PrecZeroChart{\sigma}\bigl(\PrecKernel{\sigma}\bigr).
    \]
    There exists \(\kappa_{\sigma}^{\mathrm{eqn}}>0\) such that for every
    \[
    w\in\ForSpace{\sigma},
    \qquad
    \delta p\in T_{\PrecEmbed{\sigma}(w)}\PrecSupportShell{\sigma}
    \]
    satisfying
    \[
    D\PrecData{\sigma}\,\delta p=0,
    \]
    if
    \[
    D\PrecEqChart{\sigma}(\delta p)=\delta z+\delta\eta,
    \qquad
    \delta z\in\EqnZero{\sigma},
    \quad
    \delta\eta\in\EqnHidden{\sigma},
    \]
    then
    \begin{equation}
    \begin{aligned}
    \|D\PrecShellMap{\sigma}\,\delta p\|^{2}
    &+
    \EqnPair{\sigma}\bigl(
    \PrecCharge{\sigma}\bigl((\HiddenChart{\sigma})^{-1}\delta\eta\bigr),\\
    &\qquad
    \PrecCharge{\sigma}\bigl((\HiddenChart{\sigma})^{-1}\delta\eta\bigr)
    \bigr)
    \\
    &\ge
    \kappa_{\sigma}^{\mathrm{eqn}}
    \left(
    \|\PrecStateComp{\sigma}(\delta z)\|^{2}
    +
    \|\delta\eta\|^{2}
    \right).
    \end{aligned}
    \label{eq:preccoercive}
    \end{equation}
\end{enumerate}
\end{definition}

\begin{remark}
Definition \ref{def:prec} is the precise package that now has to be explained one layer deeper. The benchmark claim boundary says that such an explanation may sharpen the origin-side derivational line, but it may not change the one operative metric, enlarge the low-energy matter law, or blur the distinction between exact relativistic adoption and first-principles derivation.
\end{remark}

\section{Precursor-Generating Substrate Origin Dynamics}

\begin{definition}[Precursor-generating substrate origin dynamics]
\label{def:orig}
Fix \(\sigma\in\{\Car,\Cur,\Hom\}\). A \emph{precursor-generating substrate origin dynamics package} for sector \(\sigma\) consists of the following data.
\begin{enumerate}
    \item A finite-dimensional Euclidean origin state space
    \[
    \OrigSpace{\sigma},
    \]
    a compact convex origin body
    \[
    \OrigBody{\sigma}\subset\OrigSpace{\sigma}
    \]
    with nonempty interior, an affine surjection
    \[
    \OrigProj{\sigma}:\OrigSpace{\sigma}\to\PrecSpace{\sigma},
    \]
    and the projected precursor body
    \[
    \PrecBody{\sigma}
    \equiv
    \OrigProj{\sigma}\bigl(\OrigBody{\sigma}\bigr)
    \subset\PrecSpace{\sigma},
    \]
    assumed compact convex with nonempty interior. There is an affine surjection
    \[
    \PrecProj{\sigma}:\PrecSpace{\sigma}\to\LiftSpace{\sigma},
    \]
    and the projected lift body
    \[
    \LiftBody{\sigma}
    \equiv
    \PrecProj{\sigma}\bigl(\PrecBody{\sigma}\bigr)
    \subset\LiftSpace{\sigma},
    \]
    also assumed compact convex with nonempty interior. There is an affine surjection
    \[
    \LiftProj{\sigma}:\LiftSpace{\sigma}\to\SubSpace{\sigma},
    \]
    and the projected substrate body
    \[
    \SubBody{\sigma}
    \equiv
    \LiftProj{\sigma}\bigl(\LiftBody{\sigma}\bigr)
    \subset\SubSpace{\sigma},
    \]
    again assumed compact convex with nonempty interior.

    There is a smooth origin precursor-fiber minimizer
    \[
    \OrigPrecMin{\sigma}:\PrecBody{\sigma}\to\OrigBody{\sigma}
    \]
    satisfying
    \[
    \OrigProj{\sigma}\circ\OrigPrecMin{\sigma}
    =
    \mathrm{id}_{\PrecBody{\sigma}},
    \]
    a smooth origin lift-fiber minimizer
    \[
    \OrigLiftMin{\sigma}:\LiftBody{\sigma}\to\OrigBody{\sigma}
    \]
    satisfying
    \[
    \PrecProj{\sigma}\circ\OrigProj{\sigma}\circ\OrigLiftMin{\sigma}
    =
    \mathrm{id}_{\LiftBody{\sigma}},
    \]
    and a smooth origin substrate-fiber minimizer
    \[
    \OrigSubMin{\sigma}:\SubBody{\sigma}\to\OrigBody{\sigma}
    \]
    satisfying
    \[
    \LiftProj{\sigma}\circ\PrecProj{\sigma}\circ\OrigProj{\sigma}\circ\OrigSubMin{\sigma}
    =
    \mathrm{id}_{\SubBody{\sigma}}.
    \]
    Define
    \begin{equation}
    \PrecLiftMin{\sigma}
    \equiv
    \OrigProj{\sigma}\circ\OrigLiftMin{\sigma},
    \qquad
    \PrecSubMin{\sigma}
    \equiv
    \OrigProj{\sigma}\circ\OrigSubMin{\sigma},
    \qquad
    \LiftSelect{\sigma}
    \equiv
    \PrecProj{\sigma}\circ\PrecSubMin{\sigma}.
    \label{eq:originducedmin}
    \end{equation}
    Assume
    \begin{equation}
    \OrigPrecMin{\sigma}\bigl(\PrecLiftMin{\sigma}(\ell)\bigr)
    =
    \OrigLiftMin{\sigma}(\ell)
    \qquad
    \text{for every }\ell\in\LiftBody{\sigma},
    \label{eq:origcompatpreclift}
    \end{equation}
    and
    \begin{equation}
    \OrigLiftMin{\sigma}\bigl(\LiftSelect{\sigma}(m)\bigr)
    =
    \OrigSubMin{\sigma}(m)
    \qquad
    \text{for every }m\in\SubBody{\sigma}.
    \label{eq:origcompatliftsub}
    \end{equation}
    \item The same finite-dimensional Euclidean substrate restriction target
    \[
    \SubTarget{\sigma},
    \]
    the same positive-definite symmetric substrate pairing
    \[
    \SubPair{\sigma}:
    \SubTarget{\sigma}\times\SubTarget{\sigma}\to\mathbb R,
    \]
    and a smooth origin restriction section
    \[
    \OrigSection{\sigma}:\OrigSpace{\sigma}\to\SubTarget{\sigma}.
    \]
    Define the origin restriction functional
    \begin{equation}
    \OrigFunctional{\sigma}(o)
    \equiv
    \SubPair{\sigma}\bigl(\OrigSection{\sigma}(o),\OrigSection{\sigma}(o)\bigr),
    \label{eq:origfunctional}
    \end{equation}
    together with the exact origin shell
    \begin{equation}
    \OrigShellSet{\sigma}
    \equiv
    \bigl(\OrigSection{\sigma}\bigr)^{-1}(0)\cap\OrigBody{\sigma}.
    \label{eq:origshell}
    \end{equation}
    Define the induced precursor restriction section
    \begin{equation}
    \PrecSection{\sigma}
    \equiv
    \OrigSection{\sigma}\circ\OrigPrecMin{\sigma},
    \label{eq:originducedprecsection}
    \end{equation}
    and the induced precursor restriction functional
    \begin{equation}
    \PrecFunctional{\sigma}(p)
    \equiv
    \SubPair{\sigma}\bigl(\PrecSection{\sigma}(p),\PrecSection{\sigma}(p)\bigr)
    =
    \OrigFunctional{\sigma}\bigl(\OrigPrecMin{\sigma}(p)\bigr).
    \label{eq:originducedprecfunctional}
    \end{equation}
    Define the induced lifted restriction section
    \begin{equation}
    \LiftSection{\sigma}
    \equiv
    \PrecSection{\sigma}\circ\PrecLiftMin{\sigma}
    =
    \OrigSection{\sigma}\circ\OrigLiftMin{\sigma},
    \label{eq:originducedliftsection}
    \end{equation}
    and the induced lifted restriction functional
    \begin{equation}
    \LiftFunctional{\sigma}(\ell)
    \equiv
    \SubPair{\sigma}\bigl(\LiftSection{\sigma}(\ell),\LiftSection{\sigma}(\ell)\bigr)
    =
    \PrecFunctional{\sigma}\bigl(\PrecLiftMin{\sigma}(\ell)\bigr)
    =
    \OrigFunctional{\sigma}\bigl(\OrigLiftMin{\sigma}(\ell)\bigr).
    \label{eq:originducedliftfunctional}
    \end{equation}
    Assume that, for each \(p\in\PrecBody{\sigma}\), \(\OrigPrecMin{\sigma}(p)\) is the unique minimizer of \(\OrigFunctional{\sigma}\) on the fiber
    \[
    \bigl(\OrigProj{\sigma}\bigr)^{-1}(p)\cap\OrigBody{\sigma},
    \]
    that, for each \(\ell\in\LiftBody{\sigma}\), \(\OrigLiftMin{\sigma}(\ell)\) is the unique minimizer of \(\OrigFunctional{\sigma}\) on the composite fiber
    \[
    \bigl(\PrecProj{\sigma}\circ\OrigProj{\sigma}\bigr)^{-1}(\ell)\cap\OrigBody{\sigma},
    \]
    and that, for each \(m\in\SubBody{\sigma}\), \(\OrigSubMin{\sigma}(m)\) is the unique minimizer of \(\OrigFunctional{\sigma}\) on the deeper composite fiber
    \[
    \bigl(\LiftProj{\sigma}\circ\PrecProj{\sigma}\circ\OrigProj{\sigma}\bigr)^{-1}(m)\cap\OrigBody{\sigma}.
    \]
    Assume that \(\PrecLiftMin{\sigma}(\ell)\) is therefore the unique minimizer of \(\PrecFunctional{\sigma}\) on the fiber
    \[
    \bigl(\PrecProj{\sigma}\bigr)^{-1}(\ell)\cap\PrecBody{\sigma},
    \]
    that \(\PrecSubMin{\sigma}(m)\) is therefore the unique minimizer of \(\PrecFunctional{\sigma}\) on the composite fiber
    \[
    \bigl(\LiftProj{\sigma}\circ\PrecProj{\sigma}\bigr)^{-1}(m)\cap\PrecBody{\sigma},
    \]
    and that \(\LiftSelect{\sigma}(m)\) is therefore the unique minimizer of \(\LiftFunctional{\sigma}\) on
    \[
    \bigl(\LiftProj{\sigma}\bigr)^{-1}(m)\cap\LiftBody{\sigma}.
    \]
    Assume there exists \(\mu_{\sigma}^{\mathrm{fib}}>0\) such that
    \begin{equation}
    \LiftFunctional{\sigma}(\ell)
    \ge
    \LiftFunctional{\sigma}\bigl(\LiftSelect{\sigma}(m)\bigr)
    +
    \mu_{\sigma}^{\mathrm{fib}}\,
    \operatorname{dist}\bigl(\ell,\LiftSelect{\sigma}(m)\bigr)^{2}
    \label{eq:origprecfibergap}
    \end{equation}
    for every \(m\in\SubBody{\sigma}\) and every \(\ell\in\bigl(\LiftProj{\sigma}\bigr)^{-1}(m)\cap\LiftBody{\sigma}\). Assume there exists \(\mu_{\sigma}^{\mathrm{sub}}>0\) such that
    \begin{equation}
    \LiftFunctional{\sigma}\bigl(\LiftSelect{\sigma}(m)\bigr)
    \ge
    \mu_{\sigma}^{\mathrm{sub}}\,
    \operatorname{dist}\bigl(m,\LiftProj{\sigma}(\PrecProj{\sigma}(\OrigProj{\sigma}(\OrigShellSet{\sigma})))\bigr)^{2}
    \qquad
    \text{for all }m\in\SubBody{\sigma}.
    \label{eq:origprecprojectedgap}
    \end{equation}
    \item Affine precursor observer-data and shell-response readouts
    \[
    \PrecData{\sigma}:\PrecSpace{\sigma}\to\DataSpace{\sigma},
    \qquad
    \PrecShellMap{\sigma}:\PrecSpace{\sigma}\to\AmbientTarget{\sigma},
    \]
    and affine origin observer-data and shell-response readouts
    \[
    \OrigData{\sigma}:\OrigSpace{\sigma}\to\DataSpace{\sigma},
    \qquad
    \OrigShellMap{\sigma}:\OrigSpace{\sigma}\to\AmbientTarget{\sigma},
    \]
    satisfying
    \begin{equation}
    \OrigData{\sigma}
    =
    \PrecData{\sigma}\circ\OrigProj{\sigma},
    \qquad
    \OrigShellMap{\sigma}
    =
    \PrecShellMap{\sigma}\circ\OrigProj{\sigma}.
    \label{eq:origfactor}
    \end{equation}
    There is also a finite-dimensional Euclidean support target
    \[
    \ResTarget{\sigma},
    \]
    a positive-definite symmetric support pairing
    \[
    \ResPair{\sigma}:
    \ResTarget{\sigma}\times\ResTarget{\sigma}\to\mathbb R,
    \]
    a finite-dimensional real hidden target
    \[
    \EqnTarget{\sigma},
    \]
    and a positive-definite symmetric hidden pairing
    \[
    \EqnPair{\sigma}:
    \EqnTarget{\sigma}\times\EqnTarget{\sigma}\to\mathbb R.
    \]
    There is a compact symmetry group
    \[
    \OrigSym{\sigma}
    \]
    and an involution
    \[
    \OrigTime{\sigma},
    \]
    acting orthogonally on \(\OrigSpace{\sigma}\), \(\PrecSpace{\sigma}\), \(\LiftSpace{\sigma}\), \(\SubSpace{\sigma}\), \(\DataSpace{\sigma}\), \(\AmbientTarget{\sigma}\), and \(\ResTarget{\sigma}\), and by \(\EqnPair{\sigma}\)-isometries on \(\EqnTarget{\sigma}\), preserving \(\OrigBody{\sigma}\), \(\PrecBody{\sigma}\), \(\LiftBody{\sigma}\), \(\OrigProj{\sigma}\), \(\PrecProj{\sigma}\), \(\LiftProj{\sigma}\), \(\OrigPrecMin{\sigma}\), \(\OrigLiftMin{\sigma}\), \(\OrigSubMin{\sigma}\), \(\OrigSection{\sigma}\), \(\PrecData{\sigma}\), \(\PrecShellMap{\sigma}\), \(\OrigData{\sigma}\), and \(\OrigShellMap{\sigma}\).
    \item The restriction of \(\OrigProj{\sigma}\) to \(\OrigShellSet{\sigma}\) is a diffeomorphism onto its image
    \[
    \PrecShellSet{\sigma}
    \equiv
    \OrigProj{\sigma}\bigl(\OrigShellSet{\sigma}\bigr)
    \subset\PrecBody{\sigma}.
    \]
    There is a Euclidean origin shell chart
    \[
    \OrigChart{\sigma}:\OrigShellSet{\sigma}\to\ResSpace{\sigma}
    \]
    whose shell-body image
    \[
    \ResBody{\sigma}
    \equiv
    \OrigChart{\sigma}\bigl(\OrigShellSet{\sigma}\bigr)
    \]
    is compact convex with nonempty interior. There exist affine maps
    \[
    \ResData{\sigma}:\ResSpace{\sigma}\to\DataSpace{\sigma},
    \qquad
    \ResShellMap{\sigma}:\ResSpace{\sigma}\to\AmbientTarget{\sigma},
    \]
    and a smooth support section
    \[
    \ResSection{\sigma}:\ResSpace{\sigma}\to\ResTarget{\sigma}
    \]
    such that
    \begin{equation}
    \ResData{\sigma}\circ\OrigChart{\sigma}
    =
    \OrigData{\sigma}\big|_{\OrigShellSet{\sigma}},
    \qquad
    \ResShellMap{\sigma}\circ\OrigChart{\sigma}
    =
    \OrigShellMap{\sigma}\big|_{\OrigShellSet{\sigma}}.
    \label{eq:origchartreadouts}
    \end{equation}
    There is a smooth origin shell-internal support recorder
    \[
    \OrigSupport{\sigma}:\OrigShellSet{\sigma}\to\ResTarget{\sigma}
    \]
    satisfying
    \begin{equation}
    \ResSection{\sigma}\circ\OrigChart{\sigma}
    =
    \OrigSupport{\sigma}.
    \label{eq:origsupporttransport}
    \end{equation}
    Define the exact origin support shell
    \begin{equation}
    \OrigSupportShell{\sigma}
    \equiv
    \bigl(\OrigSupport{\sigma}\bigr)^{-1}(0)\cap\OrigShellSet{\sigma},
    \label{eq:origsupportshell}
    \end{equation}
    and assume
    \begin{equation}
    \ResShellSet{\sigma}
    \equiv
    \bigl(\ResSection{\sigma}\bigr)^{-1}(0)\cap\ResBody{\sigma}
    =
    \OrigChart{\sigma}\bigl(\OrigSupportShell{\sigma}\bigr).
    \label{eq:origtransportedsupportshell}
    \end{equation}
    Assume there exists \(\nu_{\sigma}^{\mathrm{sup}}>0\) such that
    \begin{equation}
    \ResPair{\sigma}\bigl(\ResSection{\sigma}(u),\ResSection{\sigma}(u)\bigr)
    \ge
    \nu_{\sigma}^{\mathrm{sup}}\,
    \operatorname{dist}\bigl(u,\ResShellSet{\sigma}\bigr)^{2}
    \qquad
    \text{for all }u\in\ResBody{\sigma}.
    \label{eq:orignormal}
    \end{equation}
    Assume the restriction of \(\OrigProj{\sigma}\) to \(\OrigSupportShell{\sigma}\) is a diffeomorphism onto its image
    \[
    \PrecSupportShell{\sigma}
    \equiv
    \OrigProj{\sigma}\bigl(\OrigSupportShell{\sigma}\bigr).
    \]
    The transported symmetry and time-reversal actions act linearly on \(\ResSpace{\sigma}\), preserving \(\ResBody{\sigma}\), \(\ResShellSet{\sigma}\), \(\ResSection{\sigma}\), \(\ResData{\sigma}\), and \(\ResShellMap{\sigma}\).
    \item A closed embedded origin support zero core
    \[
    \OrigZeroCore{\sigma}\subset\OrigSupportShell{\sigma},
    \]
    a finite-dimensional Euclidean hidden-seed space
    \[
    \SeedHidden{\sigma},
    \]
    and a compact convex hidden-seed body
    \[
    \SeedHiddenBody{\sigma}\subset\SeedHidden{\sigma}
    \]
    with
    \[
    0\in\operatorname{int}\SeedHiddenBody{\sigma}.
    \]
    There is a Euclidean origin support-shell chart
    \[
    \OrigEqChart{\sigma}:\OrigSupportShell{\sigma}\to\EqnSpace{\sigma},
    \]
    an affine chart on the origin support zero core
    \[
    \OrigZeroChart{\sigma}:\OrigZeroCore{\sigma}\to\EqnZero{\sigma},
    \]
    and a linear isometric embedding
    \[
    \HiddenChart{\sigma}:\SeedHidden{\sigma}\to\EqnHidden{\sigma}
    \]
    such that
    \[
    \EqnSpace{\sigma}
    =
    \EqnZero{\sigma}\oplus\EqnHidden{\sigma}
    \]
    is an orthogonal direct sum, the set
    \[
    \ZeroBody{\sigma}
    \equiv
    \OrigZeroChart{\sigma}\bigl(\OrigZeroCore{\sigma}\bigr)
    \]
    is compact convex with nonempty interior in \(\EqnZero{\sigma}\), the shell image
    \[
    \EqnBody{\sigma}
    \equiv
    \OrigEqChart{\sigma}\bigl(\OrigSupportShell{\sigma}\bigr)
    \]
    has nonempty interior in \(\EqnSpace{\sigma}\), and there is a diffeomorphism
    \begin{equation}
    \OrigSplit{\sigma}:
    \OrigZeroCore{\sigma}\times\SeedHiddenBody{\sigma}
    \xrightarrow{\cong}
    \OrigSupportShell{\sigma}
    \label{eq:origsplit}
    \end{equation}
    satisfying
    \begin{equation}
    \OrigEqChart{\sigma}\bigl(\OrigSplit{\sigma}(z,\eta)\bigr)
    =
    \OrigZeroChart{\sigma}(z)+\HiddenChart{\sigma}(\eta).
    \label{eq:origchartsplit}
    \end{equation}
    The origin shell readouts become affine in the origin support-shell chart: there exist affine maps
    \[
    \EqnData{\sigma}:\EqnSpace{\sigma}\to\DataSpace{\sigma},
    \qquad
    \EqnShell{\sigma}:\EqnSpace{\sigma}\to\AmbientTarget{\sigma},
    \]
    such that
    \begin{equation}
    \EqnData{\sigma}\circ\OrigEqChart{\sigma}
    =
    \OrigData{\sigma}\big|_{\OrigSupportShell{\sigma}},
    \qquad
    \EqnShell{\sigma}\circ\OrigEqChart{\sigma}
    =
    \OrigShellMap{\sigma}\big|_{\OrigSupportShell{\sigma}}.
    \label{eq:origeqnchartreadouts}
    \end{equation}
    Assume the restriction of \(\OrigProj{\sigma}\) to \(\OrigZeroCore{\sigma}\) is a diffeomorphism onto its image
    \[
    \PrecZeroCore{\sigma}
    \equiv
    \OrigProj{\sigma}\bigl(\OrigZeroCore{\sigma}\bigr).
    \]
    The transported symmetry and time-reversal actions act linearly on \(\EqnSpace{\sigma}\), preserving \(\EqnZero{\sigma}\), \(\EqnHidden{\sigma}\), and \(\EqnBody{\sigma}\).
    \item A linear hidden charge map
    \[
    \OrigCharge{\sigma}:\SeedHidden{\sigma}\to\EqnTarget{\sigma}
    \]
    such that there exists \(\lambda_{\sigma}^{\mathrm{eqn}}>0\) with
    \begin{equation}
    \EqnPair{\sigma}\bigl(\OrigCharge{\sigma}\eta,\OrigCharge{\sigma}\eta\bigr)
    \ge
    \lambda_{\sigma}^{\mathrm{eqn}}\,
    \|\eta\|^{2}
    \qquad
    \text{for all }\eta\in\SeedHidden{\sigma}.
    \label{eq:origchargegap}
    \end{equation}
    Define the origin support-shell hidden recorder by
    \begin{equation}
    \OrigResidual{\sigma}\bigl(\OrigSplit{\sigma}(z,\eta)\bigr)
    \equiv
    \OrigCharge{\sigma}(\eta).
    \label{eq:origresidual}
    \end{equation}
    \item A smooth embedding
    \[
    \jmath_{\sigma}:X_{\sigma}\hookrightarrow\VisibleAffine{\sigma}
    \]
    and a smooth observer-compatible exact origin zero-response realization
    \[
    \OrigEmbed{\sigma}:\ForSpace{\sigma}\hookrightarrow\OrigZeroCore{\sigma}
    \]
    whose image is exactly
    \[
    \OrigZeroCore{\sigma}
    \cap
    \bigl(\OrigShellMap{\sigma}\bigr)^{-1}(0),
    \]
    and such that
    \begin{equation}
    \OrigData{\sigma}\bigl(\OrigEmbed{\sigma}(w)\bigr)
    =
    \bigl(
    \jmath_{\sigma}(\Reduction{\sigma}(w)),
    \CurrentCarrier{\sigma}(w),
    \EqCarrier{\sigma}(w)
    \bigr)
    \label{eq:origcompat}
    \end{equation}
    for every \(w\in\ForSpace{\sigma}\).
    \item Fixed-data solvability on the origin support zero core:
    \begin{equation}
    \OrigData{\sigma}\bigl(\OrigZeroCore{\sigma}\bigr)
    =
    \OrigData{\sigma}\Bigl(
    \OrigZeroCore{\sigma}\cap
    \bigl(\OrigShellMap{\sigma}\bigr)^{-1}(0)
    \Bigr).
    \label{eq:origsolvability}
    \end{equation}
    \item The inert origin support-zero-core kernel
    \[
    \OrigKernel{\sigma}
    \equiv
    \ker\bigl(D\OrigData{\sigma}\big|_{\OrigZeroCore{\sigma}}\bigr)
    \cap
    \ker\bigl(D\OrigShellMap{\sigma}\big|_{\OrigZeroCore{\sigma}}\bigr)
    \]
    and the orthogonal projector
    \[
    \OrigStateComp{\sigma}:
    \EqnZero{\sigma}\to
    \bigl(\OrigKernelCh{\sigma}\bigr)^{\perp}\cap\EqnZero{\sigma},
    \qquad
    \OrigKernelCh{\sigma}
    \equiv
    D\OrigZeroChart{\sigma}\bigl(\OrigKernel{\sigma}\bigr).
    \]
    There exists \(\kappa_{\sigma}^{\mathrm{eqn}}>0\) such that for every
    \[
    w\in\ForSpace{\sigma},
    \qquad
    \delta o\in T_{\OrigEmbed{\sigma}(w)}\OrigSupportShell{\sigma}
    \]
    satisfying
    \[
    D\OrigData{\sigma}\,\delta o=0,
    \]
    if
    \[
    D\OrigEqChart{\sigma}(\delta o)=\delta z+\delta\eta,
    \qquad
    \delta z\in\EqnZero{\sigma},
    \quad
    \delta\eta\in\EqnHidden{\sigma},
    \]
    then
    \begin{equation}
    \begin{aligned}
    \|D\OrigShellMap{\sigma}\,\delta o\|^{2}
    &+
    \EqnPair{\sigma}\bigl(
    \OrigCharge{\sigma}\bigl((\HiddenChart{\sigma})^{-1}\delta\eta\bigr),\\
    &\qquad
    \OrigCharge{\sigma}\bigl((\HiddenChart{\sigma})^{-1}\delta\eta\bigr)
    \bigr)
    \\
    &\ge
    \kappa_{\sigma}^{\mathrm{eqn}}
    \left(
    \|\OrigStateComp{\sigma}(\delta z)\|^{2}
    +
    \|\delta\eta\|^{2}
    \right).
    \end{aligned}
    \label{eq:origcoercive}
    \end{equation}
\end{enumerate}
\end{definition}

\begin{remark}
Definition \ref{def:orig} is deliberately deeper than the recorded precursor package but narrower than a claim of completed substrate microphysics. The precursor layer is not discarded. Instead it is shown to arise as the projected exact-shell image of a still more primitive origin layer whose exact shell, support shell, zero-core / hidden-seed split, compatibility realization, solvability statement, and coercive control all first live at the origin level.
\end{remark}

\section{Canonical Lift-Generating Precursor Package}

\begin{proposition}[The lift-generating substrate precursor package is produced canonically]
\label{prop:prec}
Assume Definition \ref{def:orig}. Define, for each sector \(\sigma\in\{\Car,\Cur,\Hom\}\),
\begin{equation}
\PrecSection{\sigma}
\equiv
\OrigSection{\sigma}\circ\OrigPrecMin{\sigma},
\label{eq:propinducedprecsection}
\end{equation}
the precursor shell chart
\begin{equation}
\PrecChart{\sigma}
\equiv
\OrigChart{\sigma}\circ
\bigl(\OrigProj{\sigma}\big|_{\OrigShellSet{\sigma}}\bigr)^{-1},
\label{eq:propinducedprecchart}
\end{equation}
the precursor shell-internal support recorder
\begin{equation}
\PrecSupport{\sigma}
\equiv
\OrigSupport{\sigma}\circ
\bigl(\OrigProj{\sigma}\big|_{\OrigShellSet{\sigma}}\bigr)^{-1},
\label{eq:propinducedprecsupport}
\end{equation}
the projected support shell and projected support zero core
\begin{equation}
\PrecSupportShell{\sigma}
\equiv
\OrigProj{\sigma}\bigl(\OrigSupportShell{\sigma}\bigr),
\qquad
\PrecZeroCore{\sigma}
\equiv
\OrigProj{\sigma}\bigl(\OrigZeroCore{\sigma}\bigr),
\label{eq:propinducedprecshells}
\end{equation}
the support-shell and zero-core charts
\begin{equation}
\PrecEqChart{\sigma}
\equiv
\OrigEqChart{\sigma}\circ
\bigl(\OrigProj{\sigma}\big|_{\OrigSupportShell{\sigma}}\bigr)^{-1},
\qquad
\PrecZeroChart{\sigma}
\equiv
\OrigZeroChart{\sigma}\circ
\bigl(\OrigProj{\sigma}\big|_{\OrigZeroCore{\sigma}}\bigr)^{-1},
\label{eq:propinducedpreccharts}
\end{equation}
the projected split map
\begin{equation}
\PrecSplit{\sigma}(z,\eta)
\equiv
\OrigProj{\sigma}\Bigl(
\OrigSplit{\sigma}\bigl(
(\OrigProj{\sigma}\big|_{\OrigZeroCore{\sigma}})^{-1}(z),
\eta
\bigr)
\Bigr),
\label{eq:propinducedprecsplit}
\end{equation}
the hidden charge and hidden recorder
\begin{equation}
\PrecCharge{\sigma}
\equiv
\OrigCharge{\sigma},
\qquad
\PrecResidual{\sigma}\bigl(\PrecSplit{\sigma}(z,\eta)\bigr)
\equiv
\OrigCharge{\sigma}(\eta),
\label{eq:propinducedprecresidual}
\end{equation}
and the exact zero-response realization
\begin{equation}
\PrecEmbed{\sigma}
\equiv
\OrigProj{\sigma}\circ\OrigEmbed{\sigma}.
\label{eq:propinducedprecembed}
\end{equation}
Then the resulting data satisfy exactly the lift-generating substrate precursor dynamics package of Definition \ref{def:prec}.
\end{proposition}

\begin{proof}
We verify the items of Definition \ref{def:prec} in order.

For item (1), the precursor state space \(\PrecSpace{\sigma}\), the compact convex precursor body \(\PrecBody{\sigma}\), the affine surjection \(\PrecProj{\sigma}\), the projected lift body \(\LiftBody{\sigma}\), the affine surjection \(\LiftProj{\sigma}\), and the projected substrate body \(\SubBody{\sigma}\) are already part of Definition \ref{def:orig}(1). Equation \eqref{eq:originducedmin} gives the smooth maps \(\PrecLiftMin{\sigma}\), \(\PrecSubMin{\sigma}\), and \(\LiftSelect{\sigma}\). Moreover,
\[
\PrecProj{\sigma}\circ\PrecLiftMin{\sigma}
=
\PrecProj{\sigma}\circ\OrigProj{\sigma}\circ\OrigLiftMin{\sigma}
=
\mathrm{id}_{\LiftBody{\sigma}},
\]
and similarly
\[
\LiftProj{\sigma}\circ\PrecProj{\sigma}\circ\PrecSubMin{\sigma}
=
\LiftProj{\sigma}\circ\PrecProj{\sigma}\circ\OrigProj{\sigma}\circ\OrigSubMin{\sigma}
=
\mathrm{id}_{\SubBody{\sigma}}.
\]
Using \eqref{eq:origcompatliftsub}, we obtain
\[
\PrecLiftMin{\sigma}\bigl(\LiftSelect{\sigma}(m)\bigr)
=
\OrigProj{\sigma}\Bigl(
\OrigLiftMin{\sigma}\bigl(\LiftSelect{\sigma}(m)\bigr)
\Bigr)
=
\OrigProj{\sigma}\bigl(\OrigSubMin{\sigma}(m)\bigr)
=
\PrecSubMin{\sigma}(m),
\]
which is exactly \eqref{eq:compatmin}.

For item (2), equation \eqref{eq:propinducedprecsection} gives the precursor restriction section, and \eqref{eq:originducedprecfunctional} gives its associated precursor restriction functional. If \(p\in\OrigProj{\sigma}(\OrigShellSet{\sigma})\), choose \(o\in\OrigShellSet{\sigma}\) with \(\OrigProj{\sigma}(o)=p\). Then \(\OrigFunctional{\sigma}(o)=0\), so by the origin precursor-fiber minimality of Definition \ref{def:orig}(2),
\[
\PrecFunctional{\sigma}(p)
=
\OrigFunctional{\sigma}\bigl(\OrigPrecMin{\sigma}(p)\bigr)
=
0.
\]
Conversely, if \(\PrecFunctional{\sigma}(p)=0\), then
\[
\OrigFunctional{\sigma}\bigl(\OrigPrecMin{\sigma}(p)\bigr)=0,
\]
so \(\OrigPrecMin{\sigma}(p)\in\OrigShellSet{\sigma}\) and therefore \(p\in\OrigProj{\sigma}(\OrigShellSet{\sigma})\). Hence
\[
\PrecShellSet{\sigma}
=
\bigl(\PrecSection{\sigma}\bigr)^{-1}(0)\cap\PrecBody{\sigma}
=
\OrigProj{\sigma}\bigl(\OrigShellSet{\sigma}\bigr).
\]
The precursor lift-fiber minimizing property of \(\PrecLiftMin{\sigma}\), the precursor substrate-fiber minimizing property of \(\PrecSubMin{\sigma}\), the lift-fiber minimizing property of \(\LiftSelect{\sigma}\), and the estimate \eqref{eq:precfibergap} are exactly the assumptions \eqref{eq:origprecfibergap} of Definition \ref{def:orig}(2). The projected lower-shell control \eqref{eq:precprojectedgap} is exactly \eqref{eq:origprecprojectedgap}.

For item (3), the affine readouts \(\PrecData{\sigma}\) and \(\PrecShellMap{\sigma}\), the lifted readouts \(\LiftData{\sigma}\) and \(\LiftShellMap{\sigma}\), the support target, support pairing, hidden target, hidden pairing, and the symmetry/time-reversal actions on \(\PrecSpace{\sigma}\) and \(\LiftSpace{\sigma}\) are already part of Definition \ref{def:orig}(3). Those same actions give the induced \(\PrecSym{\sigma}\) and \(\PrecTime{\sigma}\). Because Definition \ref{def:orig}(3) says they preserve \(\OrigProj{\sigma}\), \(\PrecProj{\sigma}\), \(\LiftProj{\sigma}\), \(\OrigPrecMin{\sigma}\), \(\OrigLiftMin{\sigma}\), \(\OrigSubMin{\sigma}\), \(\OrigSection{\sigma}\), \(\PrecData{\sigma}\), and \(\PrecShellMap{\sigma}\), they preserve \(\PrecBody{\sigma}\), \(\LiftBody{\sigma}\), \(\PrecLiftMin{\sigma}\), \(\PrecSubMin{\sigma}\), \(\LiftSelect{\sigma}\), \(\PrecSection{\sigma}\), and the full precursor readout package required in Definition \ref{def:prec}(3).

For item (4), the restriction \(\OrigProj{\sigma}\big|_{\OrigShellSet{\sigma}}\) is a diffeomorphism onto \(\PrecShellSet{\sigma}\) by Definition \ref{def:orig}(4), so \eqref{eq:propinducedprecchart} defines the precursor shell chart. Its image is exactly
\[
\ResBody{\sigma}
=
\OrigChart{\sigma}\bigl(\OrigShellSet{\sigma}\bigr).
\]
Using \eqref{eq:origfactor} and \eqref{eq:origchartreadouts},
\[
\ResData{\sigma}\circ\PrecChart{\sigma}
=
\OrigData{\sigma}\circ
\bigl(\OrigProj{\sigma}\big|_{\OrigShellSet{\sigma}}\bigr)^{-1}
=
\PrecData{\sigma}\big|_{\PrecShellSet{\sigma}},
\]
and similarly
\[
\ResShellMap{\sigma}\circ\PrecChart{\sigma}
=
\PrecShellMap{\sigma}\big|_{\PrecShellSet{\sigma}}.
\]
Equation \eqref{eq:propinducedprecsupport} defines the precursor shell-internal support recorder, and \eqref{eq:origsupporttransport} becomes exactly \eqref{eq:precsupporttransport}. Because the restriction of \(\OrigProj{\sigma}\) to \(\OrigSupportShell{\sigma}\) is a diffeomorphism onto \(\PrecSupportShell{\sigma}\), one has
\[
\ResShellSet{\sigma}
=
\OrigChart{\sigma}\bigl(\OrigSupportShell{\sigma}\bigr)
=
\PrecChart{\sigma}\bigl(\PrecSupportShell{\sigma}\bigr),
\]
which is \eqref{eq:prectransportedsupportshell}. The normal shell stability \eqref{eq:precnormal} is exactly \eqref{eq:orignormal}. The transported symmetry and time-reversal actions on \(\ResSpace{\sigma}\) are the same ones already supplied in Definition \ref{def:orig}(4).

For item (5), the projected support zero core is \(\PrecZeroCore{\sigma}\), and the restrictions of \(\OrigProj{\sigma}\big|_{\OrigSupportShell{\sigma}}\) and \(\OrigProj{\sigma}\big|_{\OrigZeroCore{\sigma}}\) are diffeomorphisms onto \(\PrecSupportShell{\sigma}\) and \(\PrecZeroCore{\sigma}\), respectively. Therefore \eqref{eq:propinducedpreccharts} defines the required charts. The orthogonal splitting
\[
\EqnSpace{\sigma}=\EqnZero{\sigma}\oplus\EqnHidden{\sigma},
\]
the compact convex bodies \(\ZeroBody{\sigma}\) and \(\EqnBody{\sigma}\), and the isometric hidden chart \(\HiddenChart{\sigma}\) are already part of Definition \ref{def:orig}(5). The projected split map \eqref{eq:propinducedprecsplit} is a diffeomorphism because \(\OrigSplit{\sigma}\) is one and \(\OrigProj{\sigma}\big|_{\OrigSupportShell{\sigma}}\) is one. Moreover, if
\[
o_{z}
\equiv
\bigl(\OrigProj{\sigma}\big|_{\OrigZeroCore{\sigma}}\bigr)^{-1}(z)
\in
\OrigZeroCore{\sigma},
\]
then
\[
\PrecEqChart{\sigma}\bigl(\PrecSplit{\sigma}(z,\eta)\bigr)
=
\OrigEqChart{\sigma}\bigl(\OrigSplit{\sigma}(o_{z},\eta)\bigr)
=
\OrigZeroChart{\sigma}(o_{z})+\HiddenChart{\sigma}(\eta)
=
\PrecZeroChart{\sigma}(z)+\HiddenChart{\sigma}(\eta),
\]
which is exactly \eqref{eq:precchartsplit}. The equation-space readouts \(\EqnData{\sigma}\) and \(\EqnShell{\sigma}\) are precisely the affine maps already supplied in Definition \ref{def:orig}(5). Using \eqref{eq:origfactor}, equation \eqref{eq:origeqnchartreadouts} becomes \eqref{eq:preceqnchartreadouts}.

For item (6), \eqref{eq:propinducedprecresidual} defines the hidden charge and hidden recorder, and the positive hidden gap \eqref{eq:precchargegap} is exactly \eqref{eq:origchargegap}.

For item (7), \eqref{eq:propinducedprecembed} defines the exact precursor zero-response realization. Since \(\OrigEmbed{\sigma}\) lands in \(\OrigZeroCore{\sigma}\cap(\OrigShellMap{\sigma})^{-1}(0)\), its projected image lands in \(\PrecZeroCore{\sigma}\cap(\PrecShellMap{\sigma})^{-1}(0)\) by \eqref{eq:origfactor}. Because \(\OrigProj{\sigma}\big|_{\OrigZeroCore{\sigma}}\) is a diffeomorphism onto \(\PrecZeroCore{\sigma}\), the image is exactly that zero-response locus. Equation \eqref{eq:preccompat} follows immediately from \eqref{eq:origcompat}, \eqref{eq:origfactor}, and \eqref{eq:propinducedprecembed}.

For item (8), fixed-data solvability \eqref{eq:precsolvability} is the projection of \eqref{eq:origsolvability} through the zero-core diffeomorphism \(\OrigProj{\sigma}\big|_{\OrigZeroCore{\sigma}}\) together with the factorization \eqref{eq:origfactor}.

For item (9), define
\[
\PrecKernel{\sigma}
\equiv
\ker\bigl(D\PrecData{\sigma}\big|_{\PrecZeroCore{\sigma}}\bigr)
\cap
\ker\bigl(D\PrecShellMap{\sigma}\big|_{\PrecZeroCore{\sigma}}\bigr),
\qquad
\PrecKernelCh{\sigma}
\equiv
D\PrecZeroChart{\sigma}\bigl(\PrecKernel{\sigma}\bigr),
\]
and let \(\PrecStateComp{\sigma}\) be the orthogonal projector onto \((\PrecKernelCh{\sigma})^{\perp}\cap\EqnZero{\sigma}\). Take \(w\in\ForSpace{\sigma}\) and \(\delta p\in T_{\PrecEmbed{\sigma}(w)}\PrecSupportShell{\sigma}\) satisfying \(D\PrecData{\sigma}\,\delta p=0\). Set
\[
\delta o
\equiv
D\bigl(\OrigProj{\sigma}\big|_{\OrigSupportShell{\sigma}}\bigr)^{-1}(\delta p)
\in
T_{\OrigEmbed{\sigma}(w)}\OrigSupportShell{\sigma}.
\]
Then \(D\OrigData{\sigma}\,\delta o=0\) by \eqref{eq:origfactor}. Writing
\[
D\PrecEqChart{\sigma}(\delta p)=\delta z+\delta\eta,
\qquad
\delta z\in\EqnZero{\sigma},
\quad
\delta\eta\in\EqnHidden{\sigma},
\]
equation \eqref{eq:origcoercive} becomes exactly \eqref{eq:preccoercive}. This completes the verification of Definition \ref{def:prec}.
\end{proof}

\begin{remark}
Proposition \ref{prop:prec} gives more than another same-output relay. Every constituent of the lift-generating substrate precursor package is now an explicit projected or transported image of deeper origin data. The precursor layer has therefore lost its status as a free insertion on the active line.
\end{remark}

\section{Benchmark Preservation}

\begin{theorem}[The same benchmark one-metric package is preserved]
\label{thm:outputs}
Assume Definition \ref{def:fixed} and, for each sector \(\sigma\in\{\Car,\Cur,\Hom\}\), precursor-generating substrate origin dynamics in the sense of Definition \ref{def:orig}. Then the benchmark output statements of \eqref{eq:fixedoutputs} remain unchanged:
\[
\MetricOut(\mathbf w)=\CompletedMetric,
\qquad
\MatterOut(\mathbf w,\psi)
=
T_{\mu\nu}^{\mathrm{matter}}[\CompletedMetric,\psi]
\]
for every \(\mathbf w\in\PrimShell\) and every matter configuration \(\psi\in\Psi\). Consequently the deeper origin theorem introduces neither a second operative metric nor an enlarged low-energy matter law.
\end{theorem}

\begin{proof}
The present note changes only the deeper origin-side status of the lift-generating substrate precursor package. It does not alter the fixed primitive shell \(\PrimShell\), the recorded metric map \(\MetricOut\), or the recorded matter map \(\MatterOut\). Those remain the continuation data of Definition \ref{def:fixed}, so \eqref{eq:fixedoutputs} continues to hold exactly.
\end{proof}

\section{Main Theorem}

\begin{theorem}[The precursor burden moves below the recorded precursor package]
\label{thm:main}
Assume the fixed primitive continuation data of Definition \ref{def:fixed} and, for each sector \(\sigma\in\{\Car,\Cur,\Hom\}\), precursor-generating substrate origin dynamics in the sense of Definition \ref{def:orig}. Then:
\begin{enumerate}
    \item the lift-generating substrate precursor dynamics are produced unchanged by projected exact-shell transport from the deeper origin dynamics;
    \item because the induced precursor package is exactly the one already used by the immediately preceding lift-generation theorem, the shell-generating substrate lift dynamics and the previously recorded downstream shell-restriction and reservoir-facing consequences remain available unchanged;
    \item the same benchmark one-metric package remains fixed: the operative metric is still exactly \(\CompletedMetric\), the ordinary matter route is unchanged, there is no second operative metric, and the low-energy matter law is not enlarged;
    \item consequently the remaining honest burden on the active line is deeper again: derive or justify the precursor-generating substrate origin dynamics themselves from still more primitive substrate structure, or else prove a sharper inevitability theorem at that origin level.
\end{enumerate}
\end{theorem}

\begin{proof}
Item (1) is Proposition \ref{prop:prec}. Item (2) is the routing consequence of item (1): the present note reproduces exactly the precursor package that the immediately preceding lift-generation theorem already uses as its input, so that theorem and its downstream consequences apply unchanged. Item (3) is Theorem \ref{thm:outputs}. Item (4) is the project-level consequence of the previous items.
\end{proof}

\begin{corollary}[What remains open after this theorem]
\label{cor:next}
After Theorem \ref{thm:main}, the remaining honest burden is no longer to justify the lift-generating substrate precursor dynamics themselves. Those dynamics are now the canonical projected exact-shell output of a sharper origin class. What remains open is why the precursor-generating substrate origin dynamics should hold, or whether an even sharper inevitability theorem can remove that still deeper origin-level freedom while preserving the same benchmark one-metric package and claim boundary.
\end{corollary}

\begin{proof}
Theorem \ref{thm:main} removes the recorded precursor package as the current unexplained stopping point. The next honest burden therefore lies strictly below it.
\end{proof}

\section{Conclusion}

The positive result of this note is narrower than a completed first-principles derivation and stronger than another precursor-level restatement. The lift-generating substrate precursor package is no longer left on the active line as a merely recorded deeper assumption. Starting from precursor-generating substrate origin dynamics, the note derives the precursor lift-fiber minimizer, the precursor substrate-fiber minimizer, the exact precursor shell as the projected origin shell, the precursor shell chart to the same reservoir body, the precursor shell-internal support recorder, the precursor support shell and precursor support zero core, the same zero-core / hidden-seed decomposition, the same hidden charge, the exact precursor zero-response realization, the fixed-data solvability statement, and the same coercive control.

Because that induced precursor package is exactly the one already used by the immediately preceding lift-generation theorem, the gain is more than another same-output origin relay. The deeper origin theorem removes the precursor layer itself as the current unexplained insertion while leaving the lift-generation theorem and its downstream shell-restriction and reservoir-facing consequences available unchanged.

The scope remains disciplined. The benchmark package is unchanged: the one operative metric is still exactly \(\CompletedMetric\), the ordinary matter route is unchanged, there is no second operative metric, and the low-energy matter law is not enlarged. This is therefore an origin-side theorem inside the active derivational program of \TheoryName{}, not a basis for stronger promotion language. The next honest burden is deeper again: derive or justify the precursor-generating substrate origin dynamics themselves from still more primitive substrate structure, or else prove a sharper inevitability theorem there. Until then, adoption and derivation should continue to be kept sharply separated.

\end{document}
